\section{L1 Introduction \& Task Decomposition}

\subsection{Vocabulary \& Platform Model}

\mult{2}

\begin{definition}{Compute Hierarchy and Vocabulary}\\
    Host \emph{delegates to} a Compute Device.
    \begin{itemize}
        \item \textbf{Host}: controller that delegates
        \item \textbf{Device}: receives the work
        \item \textbf{CU} $\to$ \textbf{PE} $\to$ \textbf{FU} (units inside the device)
        \item \textbf{UE}: process/thread running a task
    \end{itemize}
     %TODO: insert image slide 12 with notes
\end{definition}

\begin{definition}{Task}
    \begin{itemize}
        \item OS sense: a process \emph{or} a thread
        \item parallel sense: a piece of a parallel program (a unit of work)
    \end{itemize}
\end{definition}

\multend

\begin{KR}{Recipe: Map Platform $\to$ Terminology}\\
    \textbf{Find the host} which side delegates work.

    \textbf{Find the device} which side receives it; inside it CU $\to$ PE $\to$ FU.
    
    \textbf{Identify UEs} each spawned process/thread; a \emph{task} is one work slice it runs.
\end{KR}

\begin{example2}{CB1 / Jetson terminology}\\
    Programming on the CB1, delegating to the Jetson Nano. What is what?
    \tcblower
    CB1 = \important{Host}; Nano = \important{Device}; Nano SMs = CUs; CUDA cores = PEs; ALU/FPU = FU; 
    spawned threads = UEs; one Sobel slice = task.
    %TODO: insert translation from markdown
\end{example2}

\subsection{Parallel Patterns}

\begin{concept}{Pattern Hierarchy}\\
    Worked top$\to$bottom:
    \begin{enumerate}
        \item \textbf{Finding Concurrency} -- expose parallelism
        \item \textbf{Algorithm Structure} -- organise tasks
        \item \textbf{Supporting Structure} -- code/data structures
        \item \textbf{Implementation Mechanisms} -- threads, sync
    \end{enumerate}
\end{concept}

\begin{concept}{Finding Concurrency}
    Two groups of patterns:
    \begin{itemize}
        \item \textbf{Decomposition}: split by \emph{function} (task) or by \emph{data/loop}
        \item \textbf{Dependency Analysis}: Group / Order / Data-sharing
    \end{itemize}
\end{concept}

\begin{example2}{Ex.\ S22/23/25 -- Name the patterns}\\
    State the hierarchy and the Finding-Concurrency patterns.
    \tcblower
    Finding Concurrency $\to$ Algorithm Structure $\to$ Supporting Structure $\to$ Implementation Mechanisms; Finding Concurrency = \important{Decomposition} (task/data) + \important{Dependency Analysis} (group/order/share).
\end{example2}

\subsection{The Sobel Use Case}

\begin{concept}{Sobel Pipeline (6 tasks)}\\
    $\texttt{read\_image}\to\texttt{rgb\_2\_grey}\to\{\texttt{sobel\_x}\parallel\texttt{sobel\_y}\}\to\texttt{sobel\_all}\to\texttt{display\_image}$
    \begin{itemize}
        \item only \texttt{sobel\_x}\,\&\,\texttt{sobel\_y} are independent
    \end{itemize}
    % INSERT IMAGE: L1 slide 28 -- Sobel control-flow graph
\end{concept}

\begin{formula}{Sobel Operator}\\
    $G_x=\begin{bsmallmatrix}-1&0&+1\\-2&0&+2\\-1&0&+1\end{bsmallmatrix}$,\;
    $G_y=\begin{bsmallmatrix}-1&-2&-1\\0&0&0\\+1&+2&+1\end{bsmallmatrix}$
    \vspace{1mm}\\
    $|G|=\sqrt{G_x^2+G_y^2}\;\to\;|G_x|+|G_y|$ (simplified)
\end{formula}

\begin{example2}{Ex.\ S18 -- How does Sobel work?}\\
    Determine the edge-detection computation.
    \tcblower
    Two $3\times3$ convolutions give the gradient $G_x,G_y$; magnitude $|G|$ \important{simplified} to $|G_x|+|G_y|$, then threshold.
\end{example2}

\subsection{Task Decomposition}

\begin{concept}{Functional vs Loop Splitting}
    \begin{itemize}
        \item \textbf{Functional}: split by named function (few tasks, limited)
        \item \textbf{Loop/data}: split per-element loop by region (scales -- preferred)
    \end{itemize}
\end{concept}

\begin{definition}{CFG / DFG}\\
    \textbf{Control-Flow Graph}: all execution paths (nodes = actions, edges = order). \textbf{Data-Flow Graph}: data dependencies. Used to spot concurrency.
\end{definition}

\begin{KR}{Recipe: Task-Decompose an Algorithm}\\
    \textbf{Draw the CFG} nodes = functions, edges = data deps.

    \textbf{Count distinct actions} = functional tasks.

    \textbf{Find independent tasks} no path between $\Rightarrow$ concurrent.

    \textbf{Try both splittings} functional vs loop/region.

    \textbf{State parallelism + joins}
\end{KR}

\begin{example2}{Ex.\ S26--29 -- Decompose Sobel / read CFG}\\
    Split Sobel into tasks. Distinct actions? Independent tasks? Parallelisation potential?
    \tcblower
    6 functional tasks; \important{6} distinct actions; only \texttt{sobel\_x},\texttt{sobel\_y} independent. Functional parallelism limited (2-way); \important{loop/region splitting} scales.
\end{example2}

\subsection{Dependency Analysis}

\mult{2}

\begin{concept}{Group / Order / Share}
    \begin{itemize}
        \item \textbf{Group}: tasks with same data/constraints
        \item \textbf{Order}: producer before consumer
        \item \textbf{Data-sharing}: classify each datum
    \end{itemize}
\end{concept}

\begin{definition}{Data-Sharing Types}
    \begin{itemize}
        \item read-only
        \item effectively-local
        \item read-write (accumulate / multi-read-single-write)
    \end{itemize}
\end{definition}

\multend

\begin{KR}{Recipe: Dependency Analysis}
    \textbf{Group} same data/constraints together.
    \textbf{Order} add constraint per producer$\to$consumer (+ ghost-exchange before boundary).
    \textbf{Classify} read-only / local / read-write; strict ordering removes ``potentially''.
\end{KR}

\begin{example2}{Ex.\ S30--33 -- Apply to Sobel}\\
    Which sub-pattern matters most? Group/order/share for Sobel?
    \tcblower
    \important{Data Sharing}. Group \texttt{sobel\_x/y} (same grey input); order \texttt{sobel\_all} after both; grey = read-only, kernel = effectively-local, \texttt{Gx/Gy} = read-write, \texttt{sobel\_all} = \important{accumulate}.
\end{example2}

\subsection{Task Parallelism \& Fork--Join}

\begin{concept}{Task-Parallelism Checks}
    \begin{itemize}
        \item compute $>$ management overhead
        \item \#tasks $\geq$ \#PEs (more is better)
        \item minimise data dependencies
        \item equal work $\to$ static; uneven $\to$ dynamic / work-stealing
    \end{itemize}
\end{concept}

\begin{definition}{Fork--Join}
    Clone code into concurrent tasks, then \emph{join} (wait). \important{Not} POSIX \texttt{fork()}; suits recursive/irregular work.
\end{definition}

\begin{KR}{Recipe: Evaluate a Parallel Design}\\
    Check compute-vs-overhead, \#tasks-vs-\#PEs, dependencies, schedule choice.
\end{KR}

\begin{example2}{Ex.\ S34--39 -- Evaluate \& fork-join}\\
    Compute vs management? Tasks vs PEs? Is fork-join suitable for Sobel?
    \tcblower
    Whole-function task justifies a UE (per-pixel doesn't); want \#tasks $\geq$ \#PEs $\to$ region-split; only dep = \texttt{sobel\_all} join; static schedule. Fork-join: \important{poor fit} for a fixed pipeline (it suits irregular/recursive work).
\end{example2}

\subsection{Processes \& Threads}

\mult{2}

\begin{definition}{Process}
    Running program.
    \begin{itemize}
        \item unit of resource + scheduling
        \item own \texttt{main()} / address space
        \item MPU-isolated; \texttt{fork()} = CoW clone
    \end{itemize}
\end{definition}

\begin{definition}{Thread}
    Inside a process.
    \begin{itemize}
        \item shares the process's memory
        \item cheap $\to$ reuse via \important{thread pool}
    \end{itemize}
\end{definition}

\multend

\begin{example2}{Ex.\ S41--45 -- Process vs thread}\\
    What is a process? How do they differ from threads?
    \tcblower
    Process = isolated running program (own address space, \texttt{fork()}/\texttt{wait()}). Thread = lightweight, shares memory, created with \texttt{pthread\_create}; expensive to create $\to$ pool.
\end{example2}
